%!TEX root = ../template.tex
%%%%%%%%%%%%%%%%%%%%%%%%%%%%%%%%%%%%%%%%%%%%%%%%%%%%%%%%%%%%%%%%%%%
%% chapter1.tex
%% NOVA thesis document file
%%
%% Chapter with introduction
%%%%%%%%%%%%%%%%%%%%%%%%%%%%%%%%%%%%%%%%%%%%%%%%%%%%%%%%%%%%%%%%%%%

\typeout{NT FILE chapter1.tex}%

\chapter{Introduction}
\label{cha:introduction}

\section{Motivation}
\label{sec:motivation}

\ntindex[Motivation]{}

Modern computing is inherently dependent on large-scale distributed systems, where high availability and low latency are essential requirements. To meet these demands in the presence of network failures and temporary disconnections, the CAP theorem~\cite{cap-theorem} implies that strong consistency must often be relaxed in favor of availability. In this context, Conflict-Free Replicated Data Types (CRDTs) have emerged as a fundamental solution, providing Strong Eventual Consistency by construction~\cite{shapiro2011comprehensive}. CRDTs allow distributed replicas to be updated independently and concurrently, while guaranteeing that all replicas eventually converge to the same state without requiring synchronous coordination.

Despite these guarantees, implementing CRDTs correctly is challenging and prone to errors that can compromise data integrity. Consequently, formal verification techniques play an important role in increasing confidence in their correctness. Existing approaches can be divided into fully-automated verification and deductive verification. Automated tools based on SMT solvers, such as VeriFx~\cite{porre2023masses}, support validation of convergence properties and the automatic discovery of counterexamples. Deductive verification platforms, such as Why3~\cite{why3}, offer stronger guarantees by enabling proofs of functional correctness through contracts and invariants.

Experience shows that neither approach is sufficient when used in isolation. Automated verification, while efficient, struggles to handle complex or recursive structures. Deductive verification, although more expressive, often requires manual effort. This highlights the need for a methodology that combines the strengths of both approaches to support the verification of modern CRDTs.

\section{Problem Definition}
\label{sec:problem}

\ntindex[Problem Definition]{}

Despite advances in automated verification, tools such as VeriFx face clear limitations when applied to CRDTs with more complex behavior. The main problem addressed in this dissertation is that purely automated techniques are often unable to guarantee correctness in scenarios that require inductive reasoning or the preservation of global state invariants~\cite{dina2022thesis}.

Previous work has shown that SMT-based solvers exhibit well-known limitations when reasoning about recursive definitions~\cite{why3}. First, SMT-based solvers have difficulty reasoning about recursive definitions, properties that depend on iterating over data structures of unbounded size are hard to prove without explicit induction. This issue arises, for example, in vector-based counters, where computing the observable value requires recursion over a structure whose size depends on the number of replicas, often leading to timeouts or inconclusive results.

Another limitation of automated verification is that it does not guarantee that the implemented conflict-resolution policy aligns with the intended semantics. Convergence alone cannot distinguish between different behaviors, such as add-wins or remove-wins. In addition, verification results may be unstable, as small, non-functional code changes can lead to unexpected verification failures. These limitations show that automation alone is insufficient for verifying modern CRDTs, motivating the use of deductive techniques that can reason about induction, enforce invariants, and provide stronger semantic guarantees.

\section{Goals and Expected Contributions}
\label{sec:contributions}

\ntindex[Goals and Expected Contributions]{}

This work aims to explore how automated and deductive verification techniques can be combined to support the formal development of Conflict-Free Replicated Data Types. By leveraging the strengths of both approaches, it seeks to address the limitations observed when each technique is applied in isolation.

The expected contributions of this work are threefold. First, it delivers a formally verified library of CRDTs implemented in VeriFx and Why3, covering counters, sets, and register-based data types, including both state-based and operation-based models. Second, it introduces a translation methodology from VeriFx specifications to WhyML~\cite{whyml}, defining clear rules for converting convergence-oriented specifications into formal contracts and invariants suitable for deductive verification. Finally, an experimental evaluation demonstrating that the integrated approach addresses key limitations of purely automated verification, particularly in the presence of recursive definitions and global state invariants, and enables correct-by-construction reasoning in scenarios where SMT-based methods alone are insufficient.

Within this context, the specific goals of the dissertation are as follows:

\begin{itemize}
    \item To integrate distinct verification methodologies by investigating and formalizing the complementarity between VeriFx and Why3.
    \item To guarantee the correct-by-construction property by developing mechanisms that go beyond checking eventual convergence and instead enforce strong state invariants ensuring semantic validity at every execution state.
    \item To overcome expressiveness limitations in complex algorithms by addressing the verification of data structures that rely on recursive definitions or require inductive reasoning.
\end{itemize}

\section{Report Structure}
\label{sec:structure}

\ntindex[Report Structure]{}

The remainder of this document is organized as follows:

\begin{itemize}
    \item Chapter~\ref{cha:related_work} presents the state of the art on consistency models, CRDT theory, and existing formal verification tools.
    \item Chapter~\ref{cha:background} details the theoretical foundations required for this work, including the formalization of CRDTs, the logic underlying SMT solvers, and the principles of deductive verification with Why3.
    \item Chapter~\ref{cha:preliminary} describes the preliminary work carried out so far, including the implementation and verification of several CRDTs in VeriFx, the identification of practical limitations, and the successful translation and verification of selected case studies in Why3.
    \item Chapter~\ref{cha:work} outlines the work plan for completing the dissertation, including the task schedule for consolidating the proposed framework.
\end{itemize}
